\chapter{IUD (Intrauterine Device)}
\index{IUD (Intrauterine Device)}

\section*{Definition and Overview}
The intrauterine device (IUD) is a small, flexible, T-shaped device inserted into the uterus to provide long-acting, reversible contraception. It works locally in the uterus to prevent pregnancy and comes in two forms: copper IUDs, which are hormone-free and last up to 10 years, and hormonal IUDs, which release levonorgestrel and last 5 to 8 years while also reducing menstrual bleeding. The IUD is one of the most effective methods of contraception available, more than 99 per cent effective, and it is set and forget, with no daily routine. Insertion and removal are simple, quick clinic procedures.

\section*{Epidemiology}
Around one in ten women using contraception uses an IUD, and use has grown rapidly as long-acting reversible contraception has been promoted as the most effective and cost-effective approach. IUDs are suitable from adolescence through the perimenopause and are used by women of all parities. The hormonal IUD is also the most effective medical treatment for heavy menstrual bleeding and is increasingly used for this purpose. Globally, the IUD is the most widely used reversible contraceptive method.

\section*{Aetiology and Pathophysiology}
The copper IUD releases copper ions into the uterine cavity that are toxic to sperm, preventing fertilisation, and its presence alters the endometrium so that implantation cannot occur. The hormonal IUD releases a steady, low dose of progestogen that thickens cervical mucus, inhibiting sperm entry; thins the uterine lining, preventing implantation; and in many women suppresses ovulation or makes it irregular. Because the hormone acts locally, blood levels are low and systemic side effects are fewer than with the pill. The device does not interfere with the body's own hormones or fertility, and it is not an abortifacient: it prevents fertilisation and implantation.

\section*{Clinical Features}
\begin{itemize}
    \item Insertion is performed in a clinic over about five minutes
    \item Cramping and light bleeding occur during and after insertion
    \item Local anaesthetic is available to reduce discomfort
    \item The copper IUD can make periods heavier and more painful, mainly early on
    \item The hormonal IUD lightens periods and stops them in many women
    \item Spotting is common in the first months
    \item The device is checked by feeling the threads in the upper vagina
    \item A follow-up visit at three to six weeks is recommended
    \item Contraception is immediate, or from the time of fitting
\end{itemize}

\section*{Differential Diagnosis}
The decision to use an IUD is made after weighing alternatives and checking suitability.
\begin{itemize}
    \item \textbf{Contraceptive implant:} another long-acting method, placed in the arm
    \item \textbf{Pill, patch, and ring:} combined hormonal methods requiring regular use
    \item \textbf{Progestogen-only methods:} alternative for women who cannot take oestrogen
    \item \textbf{Barrier and natural methods:} less effective but non-hormonal
    \item \textbf{Sterilisation:} permanent, for those who have completed their family
    \item \textbf{Suitability issues:} pelvic infection, abnormal bleeding, and uterine anatomy
\end{itemize}

\section*{When to Seek Medical Care}
Women considering an IUD should have a full consultation, and those with risk factors for infection should be screened first. Seek review for pain, bleeding, discharge, fever, or a missed period after insertion, and for threads that cannot be felt. If you think the device has moved or fallen out, use back-up contraception and seek assessment.

\textbf{Contact your local emergency services for:}
\begin{itemize}
    \item Severe abdominal or pelvic pain
    \item Heavy bleeding with faintness, dizziness, or collapse
    \item High fever, which may indicate pelvic infection
    \item A positive pregnancy test with pain or bleeding, which may indicate ectopic pregnancy
\end{itemize}

\section*{Diagnosis and Investigations}
\begin{itemize}
    \item \textbf{Medical history and counselling:} contraception needs, medical conditions, and preferences
    \item \textbf{Pregnancy test:} to exclude pregnancy before insertion
    \item \textbf{Infection screening:} swabs and testing where indicated
    \item \textbf{Examination:} assessment of the cervix and uterus
    \item \textbf{Insertion:} placement under sterile technique, sometimes with ultrasound
    \item \textbf{Ultrasound:} if threads are missing or displacement is suspected
    \item \textbf{Follow-up:} at three to six weeks and then annually
\end{itemize}

\section*{Management}
\begin{itemize}
    \item \textbf{Insertion technique:} careful placement with the device in the fundus
    \item \textbf{Pain management:} analgesia, local anaesthetic, and reassurance
    \item \textbf{Thread checks:} teaching the woman to check threads monthly
    \item \textbf{Side effect management:} adjusting expectations and treating symptoms
    \item \textbf{Infection treatment:} antibiotics for pelvic infection if it develops
    \item \textbf{Removal and replacement:} simple procedures at the end of the device's life
    \item \textbf{Back-up contraception:} if the device is expelled or removed early
    \item \textbf{Ongoing screening:} cervical screening and STI testing as recommended
\end{itemize}

\section*{Complications}
\begin{itemize}
    \item Expulsion, occurring in around one in twenty women, usually in the first year
    \item Perforation of the uterus, rare, and usually at insertion
    \item Pelvic infection, mainly in the first three weeks
    \item Bleeding changes and cramping
    \item Missing threads, requiring assessment
    \item Pregnancy with the device in situ, which needs specialist care
    \item Rarely, ovarian cysts with the hormonal IUD, which usually resolve
\end{itemize}

\section*{Prognosis and Outcomes}
The IUD is extremely effective and safe, and satisfaction and continuation rates are among the highest of any contraceptive. Side effects settle for most women within months, and the hormonal IUD treats heavy periods in most users. Fertility returns immediately after removal, and there is no delay in conception. With a follow-up visit and simple checks, most women use the IUD successfully for its full lifespan.

\section*{Prevention and Self-Care}
\begin{itemize}
    \item Discuss all options to make an informed choice
    \item Choose an experienced provider for insertion
    \item Take pain relief before and after insertion as advised
    \item Use back-up contraception until the device is effective
    \item Check threads monthly
    \item Report pain, fever, discharge, or missed periods
    \item Attend follow-up and annual reviews
    \item Plan for replacement or removal in advance
    \item Seek early review if you think the device has moved
\end{itemize}

\section*{Sources and Further Reading}
The following reputable sources provide further information for healthcare students and patients seeking reliable, current advice about the intrauterine device.
\begin{itemize}
    \item Family Planning Australia. \textit{IUD information.} https://www.fpnsw.org.au/
    \item Healthdirect Australia. \textit{Intrauterine device.} https://www.healthdirect.gov.au/
    \item Jean Hailes for Women's Health. \textit{Contraception.} https://www.jeanhailes.org.au/
    \item Sexual Health Victoria. \textit{IUD information and services.} https://shvic.org.au/
    \item NHS. \textit{Contraception: IUD.} https://www.nhs.uk/conditions/contraception/iud-coil/
    \item World Health Organization. \textit{Family planning and contraceptives.} https://www.who.int/
\end{itemize}
